\section*{Chapitre \ref{ch:b3:banach} --- Espaces de Banach et
théorèmes fondamentaux}

\begin{solution}{exo:b3:banach:1}
(a) $\norm{Sx}_2 = \norm x_2$~: $S$ est une isométrie,
$\vertiii S = 1$. Pour le décalage arrière~: $\norm{S^*x}_2^2
= \sum_{n \geq 2}\abs{x_n}^2 \leq \norm x_2^2$, avec égalité
pour $x = e_2$~: $\vertiii{S^*} = 1$.

(b) $\norm{M_ax}_2^2 = \sum\abs{a_n}^2\abs{x_n}^2 \leq \norm
a_\infty^2\norm x_2^2$~; tester $x = e_n$ donne $\vertiii{M_a}
\geq \abs{a_n}$ pour tout $n$~: $\vertiii{M_a} = \norm
a_\infty$.

(c) $\abs{\Lambda(f)} \leq \int_0^1\abs f \leq \norm
f_\infty$~: $\norm\Lambda \leq 1$. Pour $\varepsilon > 0$ soit
$f_\varepsilon$ égale à $1$ sur $[0, \frac12 - \varepsilon]$,
$-1$ sur $[\frac12 + \varepsilon, 1]$, affine entre~:
$\norm{f_\varepsilon}_\infty = 1$ et $\Lambda(f_\varepsilon)
\geq 1 - 2\varepsilon$~: $\norm\Lambda = 1$. Non atteinte~:
$\Lambda(f) = 1$ avec $\norm f_\infty \leq 1$ force
$\int_0^{1/2}f = \frac12$ et $\int_{1/2}^1 f = -\frac12$,
c'est-à-dire (continuité, $\abs f \leq 1$) $f \equiv 1$ sur
$[0, \frac12]$ et $f \equiv -1$ sur $[\frac12, 1]$~:
contradiction en $\frac12$.
\end{solution}

\begin{solution}{exo:b3:banach:2}
Écrire $S = T\bigl(I - T^{-1}(T - S)\bigr)$ avec
$\vertiii{T^{-1}(T-S)} \leq \vertiii{T^{-1}}\,\vertiii{T - S}
= \theta < 1$~: par \cref{prop:b3:banach:neumann}, $S$ est
inversible avec $S^{-1} =
\sum_{n\geq0}\bigl(T^{-1}(T-S)\bigr)^nT^{-1}$, d'où
\[
\vertiii{S^{-1} - T^{-1}} \leq
\sum_{n\geq1}\theta^n\,\vertiii{T^{-1}}
= \frac{\vertiii{T^{-1}}^2\,\vertiii{T - S}}{1 - \theta}.
\]
Pour le système linéaire~: $x_T = T^{-1}b$ et $x_S = S^{-1}b$
diffèrent d'au plus cette borne fois $\norm b$ --- les
petites perturbations d'un système inversible restent
uniquement solvables, avec dépendance lipschitzienne de la
solution en l'opérateur.
\end{solution}

\begin{solution}{exo:b3:banach:3}
(a) $\ell^1$~: soit $(x^{(k)})$ de Cauchy. Chaque coordonnée
est de Cauchy ($\abs{x^{(k)}_n - x^{(l)}_n} \leq
\norm{x^{(k)} - x^{(l)}}_1$)~: soit $x_n =
\lim_kx^{(k)}_n$. Étant donné $\varepsilon$, pour $k, l \geq
K$~: $\sum_{n \leq N}\abs{x^{(k)}_n - x^{(l)}_n} \leq
\varepsilon$ pour tout $N$~; faire $l \to \infty$, puis $N \to
\infty$~: $\norm{x^{(k)} - x}_1 \leq \varepsilon$, et $x =
x^{(k)} - (x^{(k)} - x) \in \ell^1$. $\ell^\infty$~: Cauchy
pour $\norm\cdot_\infty$ est uniformément de Cauchy~: converge
uniformément vers une suite bornée. $c_0$ est fermé dans
$\ell^\infty$~: si $x^{(k)} \to x$ uniformément avec
$x^{(k)}_n \to_n 0$, alors $\abs{x_n} \leq \norm{x -
x^{(k)}}_\infty + \abs{x^{(k)}_n}$ donne
$\limsup_n\abs{x_n} \leq \varepsilon$~: $x \in c_0$~; un
sous-espace fermé d'un Banach est Banach. Suites finies~: leur
\omterm{def:b3:topology:interior}{adhérence} contient tout $x \in c_0$ (les troncatures
convergent~: $\sup_{n>N}\abs{x_n} \to 0$) et est contenue
dans le fermé $c_0$.

(b) Par homogénéité supposer $\norm x_p = 1$~: alors
$\abs{x_n} \leq 1$ pour tout $n$, donc $\abs{x_n}^q \leq
\abs{x_n}^p$ et $\norm x_q \leq 1 = \norm x_p$~; pour $q =
\infty$, $\abs{x_n} \leq \norm x_p$ directement. Stricte~:
$x_n = n^{-\alpha}$ avec $\frac1q < \alpha \leq \frac1p$ est
dans $\ell^q \setminus \ell^p$ (séries de Riemann).
\end{solution}

\begin{solution}{exo:b3:banach:4}
($\leq$) Si $\norm f \leq 1$ et $f\restriction_F = 0$~: pour
tout $y \in F$, $\abs{f(x)} = \abs{f(x - y)} \leq \norm{x -
y}$~; prendre l'infimum. ($\geq$, atteint)
\cref{cor:b3:banach:hbnormed}(3) produit $f$ avec
$f\restriction_F = 0$, $\norm f \leq 1$, $f(x) = d(x, F)$~: le
supremum est un maximum. Conséquence~: $F \subseteq
\bigcap\{\ker f : f\restriction_F = 0\}$ trivialement, et un
point $x \notin F$ est exclu de l'intersection par la
fonctionnelle ci-dessus ($f(x) = d(x,F) > 0$, $F$ étant
fermé).
\end{solution}

\begin{solution}{exo:b3:banach:5}
Pour $y \in \ell^\infty$~: $\abs{\Lambda_y(x)} =
\abs{\sum x_ny_n} \leq \norm y_\infty\norm x_1$, donc
$\norm{\Lambda_y} \leq \norm y_\infty$~; tester sur $e_n$~:
$\abs{y_n} = \abs{\Lambda_y(e_n)} \leq \norm{\Lambda_y}$~:
égalité. Réciproquement, étant donné $\Lambda \in
(\ell^1)'$, poser $y_n = \Lambda(e_n)$~: $\abs{y_n} \leq
\norm\Lambda$, donc $y \in \ell^\infty$~; $\Lambda$ et
$\Lambda_y$ coïncident sur les suites finies, \omterm{def:b3:topology:interior}{denses} dans
$\ell^1$ ($\norm{x - \sum_{n\leq N}x_ne_n}_1 =
\sum_{n>N}\abs{x_n} \to 0$)~: $\Lambda = \Lambda_y$.
L'application $y \mapsto \Lambda_y$ est linéaire, isométrique,
surjective. Pour $(\ell^\infty)'$ le même départ produit une
suite $y_n = \Lambda(e_n)$, mais les suites finies ne sont
\emph{pas} \omterm{def:b3:topology:interior}{denses} dans $\ell^\infty$ (la suite constante
$\mathbf 1$ est à distance $1$ de toutes), donc $\Lambda$
n'est pas déterminée par les $y_n$ --- et en effet
$(\ell^\infty)' \neq \ell^1$ (\cref{pb:b3:banach:1}).
\end{solution}

\begin{solution}{exo:b3:banach:6}
Pour chaque $x$ fixé, $y \mapsto B(x, y)$ est linéaire
\omterm{def:b3:topology:continuity}{continue}~: $\sup_{\norm y \leq 1}\norm{B(x,y)} < \infty$.
Donc la famille $\{B(\cdot, y) : \norm y \leq 1\} \subseteq
\mathcal L(E, G)$ (chaque membre \omterm{def:b3:topology:continuity}{continue}, par continuité en
$x$) est ponctuellement bornée sur le Banach $E$~:
Banach--Steinhaus (\cref{thm:b3:banach:banachsteinhaus})
fournit $C$ avec $\norm{B(x,y)} \leq C\norm x$ pour tous
$\norm y \leq 1$~; l'homogénéité en $y$ termine~:
$\norm{B(x,y)} \leq C\norm x\norm y$.
\end{solution}

\begin{solution}{exo:b3:banach:7}
(a) $\norm f_1 \leq \norm f_\infty$~: l'identité est bornée et
bijective. Son inverse est non borné~: $f_n(x) = x^n$ a
$\norm{f_n}_1 = \frac1{n+1} \to 0$ mais $\norm{f_n}_\infty =
1$. Pas de contradiction avec
\cref{cor:b3:banach:equivnorms}~: $(\mathcal C(\intcc01),
\norm\cdot_1)$ n'est \emph{pas} \omterm{def:b3:complete:complete}{complet}
(\cref{exo:b3:complete:1})~; le corollaire exige la
complétude des deux côtés.

(b) Sur l'espace $E_0$ des suites finies (\omterm{def:b3:topology:interior}{dense} dans
$\ell^2$), $\varphi(x) = \sum_n n\,x_n$ est linéaire et non
bornée ($\varphi(e_n) = n$ avec $\norm{e_n}_2 = 1$). Le
théorème du graphe fermé ne s'applique pas~: $E_0$ n'est pas
\omterm{def:b3:complete:complete}{complet} --- et $\varphi$ n'a pas d'extension \omterm{def:b3:topology:continuity}{continue} à
$\ell^2$, illustrant que la densité sans continuité uniforme
est impuissante (\cref{thm:b3:complete:extension}).
\end{solution}

\begin{solution}{exo:b3:banach:8}
On vérifie l'hypothèse du graphe fermé. Soit $x_k \to x$ et
$Tx_k \to z$ dans $\ell^2$. Pour tout $y$~:
\[
\langle z, y\rangle = \lim_k\langle Tx_k, y\rangle
= \lim_k \langle x_k, Ty\rangle = \langle x, Ty\rangle
= \langle Tx, y\rangle,
\]
en utilisant la continuité du produit scalaire dans chaque
créneau (Cauchy--Schwarz) et la symétrie deux fois. Donc $z -
Tx$ est orthogonal à tout $y$, en particulier à lui-même~: $z
= Tx$. Le graphe est fermé et $\ell^2$ est de Banach~: $T$
est borné (\cref{thm:b3:banach:closedgraph}). Donc un
opérateur symétrique défini sur \emph{tout} $\ell^2$ est
automatiquement borné~; les opérateurs symétriques
véritablement non bornés (position, impulsion, hamiltoniens)
doivent vivre sur des sous-espaces \omterm{def:b3:topology:interior}{denses} propres.
\end{solution}

\begin{solution}{exo:b3:banach:9}
D'abord, $\norm{\Lambda_n} = \sum_i\abs{w_{i,n}}$~: $\leq$ est
l'inégalité triangulaire~; $\geq$ en testant une $f$ affine
par morceaux avec $\norm f_\infty \leq 1$ et $f(x_{i,n}) =
\operatorname{sign}(w_{i,n})$ (interpoler linéairement entre
les nœuds finiment nombreux).

($\Rightarrow$) La convergence ponctuelle en toute $f$
implique (i), et la bornitude ponctuelle, donc
Banach--Steinhaus (\cref{thm:b3:banach:banachsteinhaus}) sur
le Banach $\mathcal C(\intcc01)$ donne (ii).

($\Leftarrow$) Soit $M = \sup_n\norm{\Lambda_n} + 1$. Étant
donné $f$ et $\varepsilon$, choisir un polynôme $P$ avec
$\norm{f - P}_\infty < \varepsilon/(2M)$
(\cref{cor:b3:complete:weierstrass})~; alors
\[
\Bigl|\Lambda_n f - \int_0^1 f\Bigr|
\leq \abs{\Lambda_n(f - P)} + \Bigl|\Lambda_nP -
\int_0^1P\Bigr| + \Bigl|\int_0^1(P - f)\Bigr|
\leq \varepsilon + \Bigl|\Lambda_nP - \int_0^1P\Bigr| \to
\varepsilon .
\]
Poids positifs, exactitude sur les constantes~:
$\sum_i\abs{w_{i,n}} = \sum_iw_{i,n} = \Lambda_n(\mathbf 1)
= \int_0^1 1 = 1$ pour les règles exactes sur les constantes
--- (ii) tient avec constante $1$.
\end{solution}

\begin{solution}{exo:b3:banach:10}
La linéarité de $J$ est formelle~; $\norm{J(x)} =
\sup_{\norm f \leq 1}\abs{f(x)} = \norm x$ par
\cref{cor:b3:banach:hbnormed}(2). Si $\dim E = n$~: $\dim E'
= n$ (une \omterm{def:b3:topology:basis}{base} donne des fonctionnelles coordonnées), donc
$\dim E'' = n$, et le $J$ injectif (isométrique) est
surjectif. Séparabilité~: soit $(f_n)$ \omterm{def:b3:topology:interior}{dense} dans $E'$ et
choisir $\norm{x_n} = 1$ avec $\abs{f_n(x_n)} \geq
\frac12\norm{f_n}$. Soit $F =
\overline{\operatorname{Vect}}(x_n)$~; si $F \neq E$, prendre
$g \in E'$, $g \neq 0$, s'annulant sur $F$
(\cref{cor:b3:banach:hbnormed}(3))~; choisir $f_{n_k} \to
g$~:
\[
\norm{f_{n_k} - g} \geq \abs{(f_{n_k} - g)(x_{n_k})}
= \abs{f_{n_k}(x_{n_k})} \geq \tfrac12\norm{f_{n_k}}
\geq \tfrac12\bigl(\norm g - \norm{f_{n_k} - g}\bigr),
\]
donc $\norm{f_{n_k} - g} \geq \frac13\norm g > 0$~:
contradiction. D'où $F = E$, et les combinaisons rationnelles
(ou $\Q + \iu\Q$) des $x_n$ forment un \omterm{def:b3:topology:interior}{dense} dénombrable.
\end{solution}

\begin{solution}{exo:b3:banach:11}
(a) Bien définie~: $d(x, F)$ ne dépend que de $\bar x$
(translater $x$ par $F$ ne change pas la distance).
Homogénéité et inégalité triangulaire passent de $\norm\cdot$
via l'infimum. La séparation demande la fermeture~:
$\norm{\bar x} = 0$ signifie $d(x, F) = 0$, c'est-à-dire $x
\in \bar F = F$, c'est-à-dire $\bar x = 0$. $\vertiii\pi \leq
1$~: $\norm{\bar x} \leq \norm x$. Boule ouverte sur boule
ouverte~: si $\norm{\bar x} < 1$, quelque représentant a
$\norm{x - y} < 1$~; réciproquement $\pi(B_E(0,1)) \subseteq
B_{E/F}(0,1)$ par l'inégalité de norme --- donc $\pi$ est
ouverte, le cas modèle du théorème de l'application ouverte.

(b) Soit $\sum_k\norm{\bar x_k} < \infty$ et relever avec
$\norm{x_k} \leq \norm{\bar x_k} + 2^{-k}$~: alors
$\sum\norm{x_k} < \infty$, donc $s = \sum_kx_k$ converge dans
le Banach $E$ (\cref{exo:b3:complete:1}(b)), et la continuité
de $\pi$ donne $\sum_k\bar x_k = \bar s$~: toute série
absolument convergente de $E/F$ converge, ce qui équivaut à
la complétude (même exercice).

(c) L'application $\lambda(x) = \lim_nx_n$ est linéaire $c
\to K$, s'annule exactement sur $c_0$, donc induit une
bijection linéaire $c/c_0 \to K$. Isométrie~: $d(x, c_0) =
\abs{\lambda(x)}$ --- $\leq$~: soustraire de $x$ la suite $x
- \lambda(x)\mathbf 1 \in c_0$, laissant $\lambda(x)\mathbf
1$ de norme $\abs{\lambda(x)}$~; $\geq$~: pour $y \in c_0$,
$\norm{x - y}_\infty \geq \limsup_n\abs{x_n - y_n} =
\abs{\lambda(x)}$.
\end{solution}

\begin{solution}{exo:b3:banach:12}
(a) $W = \ker P$ est fermé (image réciproque de $0$ sous une
\omterm{def:b3:topology:continuity}{application continue})~; $V = \operatorname{im}P = \ker(I -
P)$ (en effet $Px = x$ ssi $x \in \operatorname{im}P$, en
utilisant $P^2 = P$), fermé de même. Tout $x$ se scinde en
$Px + (x - Px)$ avec $Px \in V$, $x - Px \in W$, et $V \cap W
= 0$ ($x = Px = 0$)~: $E = V \oplus W$.

(b) L'argument du graphe~: soit $x_n \to x$ et $Px_n \to z$.
Alors $z \in V$ ($V$ fermé, $Px_n \in V$) et $x_n - Px_n \to
x - z \in W$ ($W$ fermé). Donc $x = z + (x - z)$ avec $z \in
V$, $x - z \in W$~; par unicité de la décomposition, $z =
Px$. Le graphe de $P$ est fermé, $E$ est de Banach~: $P$ est
bornée (\cref{thm:b3:banach:closedgraph}).

(c) (a) et (b) ensemble sont l'équivalence. Dans un espace de
Hilbert tout $V$ fermé a le complément fermé $V^\perp$
(\cref{ch:b3:hilbert})~: tout sous-espace fermé est
complémenté. Dans les Banach généraux cela échoue --- $c_0$
n'a pas de complément fermé dans $\ell^\infty$ (théorème de
Phillips, hors de nos outils) --- donc les projections
bornées sont un privilège, et le théorème du graphe fermé est
exactement la comptabilité qui convertit les scindages
géométriques en opérateurs bornés.
\end{solution}
\begin{solution}{pb:b3:banach:1}
\textbf{1.} Pour $a, b > 0$~: par concavité de $\log$,
$\log\bigl(\tfrac{a^p}p + \tfrac{b^q}q\bigr) \geq
\tfrac1p\log a^p + \tfrac1q\log b^q = \log(ab)$~;
exponentier. (Si $ab = 0$ l'inégalité est triviale.) Égalité
ssi $a^p = b^q$.

\textbf{2.} On peut supposer $\norm x_p = \norm y_q = 1$
(homogénéité). Alors
\[
\abs{\langle x, y\rangle} \leq \sum_n\abs{x_n}\abs{y_n}
\leq \sum_n\Bigl(\frac{\abs{x_n}^p}p +
\frac{\abs{y_n}^q}q\Bigr) = \frac1p + \frac1q = 1.
\]
L'égalité demande $\abs{x_n}^p = \abs{y_n}^q$ pour tout $n$ et
l'alignement des phases de $x_ny_n$.

\textbf{3.} $\abs{x_n + y_n}^p \leq \bigl(\abs{x_n} +
\abs{y_n}\bigr)\abs{x_n + y_n}^{p-1}$~; sommer et appliquer
Hölder ($p$ contre $q$, notant $(p - 1)q = p$)~:
\[
\norm{x + y}_p^p \leq \bigl(\norm x_p + \norm
y_p\bigr)\,\norm{x+y}_p^{p/q};
\]
si $\norm{x + y}_p \neq 0$, diviser par $\norm{x+y}_p^{p/q}$
et utiliser $p - \frac pq = 1$.

\textbf{4.} Complétude~: comme pour $\ell^1$
(\cref{exo:b3:banach:3}), limites coordonnée par coordonnée
plus borne uniforme de queue. Densité des suites finies~:
$\norm{x - \sum_{n\leq N}x_ne_n}_p^p = \sum_{n >
N}\abs{x_n}^p \to 0$.

\textbf{5.} Hölder donne $\norm{\Lambda_y} \leq \norm y_q$.
Tester~: $x^{(N)}_n =
\abs{y_n}^{q-1}\overline{\operatorname{sign}}(y_n)$ pour $n
\leq N$~: $\Lambda_y(x^{(N)}) = \sum_{n\leq N}\abs{y_n}^q$ et
$\norm{x^{(N)}}_p = \bigl(\sum_{n\leq
N}\abs{y_n}^{q}\bigr)^{1/p}$, d'où $\norm{\Lambda_y} \geq
\bigl(\sum_{n\leq N}\abs{y_n}^q\bigr)^{1/q} \to \norm y_q$.

\textbf{6.} Poser $y_n = \Lambda(e_n)$. Avec les mêmes
vecteurs tests, $\bigl(\sum_{n\leq N}\abs{y_n}^q\bigr)^{1/q}
\leq \norm\Lambda$ pour tout $N$~: $y \in \ell^q$. Les
fonctionnelles $\Lambda$ et $\Lambda_y$ coïncident sur les
suites finies \omterm{def:b3:topology:interior}{denses}~: $\Lambda = \Lambda_y$. Avec la
question~5, $y \mapsto \Lambda_y$ est un isomorphisme
isométrique $\ell^q \cong (\ell^p)'$.

\textbf{7.} Soit $\xi \in (\ell^p)''$. En composant avec
l'isométrie $\ell^q \cong (\ell^p)'$, $\xi$ définit un
élément de $(\ell^q)'$, qui (question~6 avec $p, q$
échangés) est $\Lambda_z$ pour un unique $z \in \ell^p$.
Alors $\xi(\Lambda_y) = \sum_nz_ny_n = J(z)(\Lambda_y)$~:
$\xi = J(z)$, $J$ est surjective --- $\ell^p$ est réflexif.

\textbf{8.} Les suites finies à entrées rationnelles sont
dénombrables et \omterm{def:b3:topology:interior}{denses} dans $\ell^p$ ($p < \infty$) et dans
$c_0$. Dans $\ell^\infty$~: la famille $\{\mathbf 1_A : A
\subseteq \N\}$ est non dénombrable avec $\norm{\mathbf 1_A -
\mathbf 1_B}_\infty = 1$ pour $A \neq B$~: pas de \omterm{def:b3:topology:interior}{dense}
dénombrable.

\textbf{9.} Si $\ell^1$ était réflexif, alors $(\ell^\infty)'
\cong J(\ell^1)$ serait \omterm{def:b3:galois:separable}{séparable}~; par
\cref{exo:b3:banach:10}, la séparabilité du \omterm{def:b3:banach:operator}{dual} forcerait
$\ell^\infty$ \omterm{def:b3:galois:separable}{séparable} --- contredisant la question~8. Donc
$\ell^1$ n'est pas réflexif.

\textbf{10.} Homogénéité de $p$ claire~; sous-additivité~:
$\limsup(u_n + v_n) \leq \limsup u_n + \limsup v_n$. Sur
$c$~: les moyennes de Cesàro d'une suite convergente
convergent vers sa limite, donc $p(x) = \lim x =
\mathrm{LIM}(x)$ là. Hahn--Banach
(\cref{thm:b3:banach:hahnbanach}) étend $\mathrm{LIM}$ à
$\ell^\infty_\R$ avec $\mathrm{LIM} \leq p$. Positivité~:
pour $x \geq 0$, $-\mathrm{LIM}(x) = \mathrm{LIM}(-x) \leq
p(-x) \leq 0$. Invariance par décalage~: les moyennes de $x -
Sx$ se télescopent en $\frac{x_1 - x_{n+1}}n \to 0$, donc
$p(\pm(x - Sx)) = 0$ et $\mathrm{LIM}(x - Sx) = 0$. Bornes~:
$\mathrm{LIM}(x) \leq p(x) \leq \limsup x$, et appliquer à
$-x$ pour la borne inférieure.

\textbf{11.} $e_n \in c_0$, donc $\mathrm{LIM}(e_n) = 0$. Si
$\mathrm{LIM} = \Lambda_y$ avec $y \in \ell^1$, alors $y_n =
0$ pour tout $n$~: $\Lambda_y = 0$~; mais
$\mathrm{LIM}(\mathbf 1) = 1$. Donc $(\ell^\infty)' \neq
\ell^1$.

\textbf{12.} Pour $x = (0,1,0,1,\dots)$~: $x + Sx = \mathbf
1$, donc $2\,\mathrm{LIM}(x) = 1$~: $\mathrm{LIM}(x) =
\frac12$. Si $\varphi$ était multiplicative et invariante par
décalage~: $x\cdot Sx = 0$ donne $\varphi(x)^2 = 0$, donc
$\varphi(x) = 0$~; mais $\varphi(x) + \varphi(Sx) = 1$ donne
$2\varphi(x) = 1$~: contradiction.

\textbf{13.} Soit $\norm{x^{(k)}}_p \leq M$. Les coordonnées
sont bornées par $M$~: une extraction diagonale donne une
sous-suite avec $x^{(k)}_n \to x_n$ pour tout $n$. Alors $x
\in \ell^p$ (bornes de sections finies). Convergence faible~:
pour $y \in \ell^q$, découper en tête finie + queue Hölder.
Dans $\ell^1$ cela échoue~: $(e_n)$ est bornée, toute
sous-suite converge coordonnée par coordonnée vers $0$, mais
tester contre $y_{n_k} = (-1)^k$ fait diverger
$\Lambda_y(e_{n_k})$.

\textbf{14.} $(c_0)' \cong \ell^1$ par le schéma des
questions~5--6 avec tests
$x^{(N)} = (\operatorname{sign}y_1, \dots,
\operatorname{sign}y_N, 0, \dots) \in c_0$. Chaîne des
duaux~: $(c_0)' = \ell^1$, $(\ell^1)' = \ell^\infty$,
$(\ell^\infty)' \supsetneq \ell^1$~: la réflexivité échoue
dès la première étape ($c_0'' = \ell^\infty \neq c_0$).

\textbf{15.} $J x^{(k)} \in E'' = (E')'$ est une famille de
fonctionnelles bornées sur le Banach $E'$~; pour chaque
$\Lambda$, $Jx^{(k)}(\Lambda) = \Lambda(x^{(k)})$ converge,
donc est bornée. Banach--Steinhaus sur $E'$ donne
$\sup_k\norm{Jx^{(k)}} < \infty$, et $J$ est isométrique~:
$\sup_k\norm{x^{(k)}} < \infty$.

\textbf{16.} ($\Rightarrow$) Bornitude~: question~15~;
coordonnées~: les $\Lambda_{e_n}$. ($\Leftarrow$) Approximer
$y \in \ell^q$ par des suites finies et utiliser la
bornitude. Pour $e_k$ dans $\ell^2$~: bornée, coordonnée par
coordonnée $\to 0$, donc $e_k \rightharpoonup 0$, pourtant
$\norm{e_k} = 1$.

\textbf{17.} Prendre $\Lambda$ avec $\norm\Lambda = 1$ et
$\Lambda(x) = \norm x$
(\cref{cor:b3:banach:hbnormed})~: $\norm x =
\lim\Lambda(x^{(k)}) \leq \liminf\norm{x^{(k)}}$.

\textbf{18.} Dans $\ell^2$, $\norm{x^{(k)} - x}_2^2 =
\norm{x^{(k)}}_2^2 - 2\operatorname{Re}\langle x^{(k)},
x\rangle + \norm x_2^2$. Convergence faible de $\Lambda_x$ et
convergence des normes~: le second membre $\to 0$.
Contre-exemple sans convergence de normes~: $e_k
\rightharpoonup 0$, $\norm{e_k} = 1 \not\to 0$.

\textbf{19.} Convergence coordonnée~: appliquer
$\Lambda_{e_n}$. Construction des blocs~: choisir $k_j$ assez
grand pour que la tête jusqu'à $N_{j-1}$ porte moins de
$\frac\delta{20}$ de masse, puis $N_j$ assez grand pour que
la queue au-delà de $N_j$ porte moins de $\frac\delta{20}$~:
le bloc $B_j$ porte tout sauf $\frac\delta{10}$.

\textbf{20.} Avec $y_n = \operatorname{sign}(x^{(k_j)}_n)$ sur
$B_j$~: $\langle x^{(k_j)}, y\rangle \geq
\norm{x^{(k_j)}}_1 - \frac{2\delta}{10} \geq
\frac{4\delta}5 > 0$, contredisant $x^{(k)}
\rightharpoonup 0$. Donc les suites faiblement nulles de
$\ell^1$ sont nulles en norme~: théorème de Schur.

\textbf{21.} Une sous-suite faiblement convergente de $(e_k)$
aurait pour limite $0$ (coordonnées), donc par Schur
$\norm{e_{k_j}}_1 \to 0$ --- mais les normes valent $1$. Pas
de paradoxe~: Schur dit que les suites ne distinguent pas les
\omterm{def:b3:topology:topology}{topologies} faible et en norme dans $\ell^1$, mais la
\omterm{thm:b3:topology:metriccompact}{compacité séquentielle} faible de la boule est une propriété
plus forte, équivalente à la réflexivité.

\textbf{22.} Tableau~: $c_0$ \omterm{def:b3:banach:operator}{dual} $\ell^1$, \omterm{def:b3:galois:separable}{séparable}, non
réflexif, pas de cpt.\ faible seq.\ ($e_1+\dots+e_k$)~;
$\ell^1$ \omterm{def:b3:banach:operator}{dual} $\ell^\infty$, \omterm{def:b3:galois:separable}{séparable}, non réflexif, pas de
cpt.\ faible seq.\ ($e_k$, Schur)~; $\ell^p$ \omterm{def:b3:banach:operator}{dual} $\ell^q$,
\omterm{def:b3:galois:separable}{séparable}, réflexif, cpt.\ faible seq.\ oui~; $\ell^\infty$
\omterm{def:b3:banach:operator}{dual} $\supsetneq\ell^1$, non \omterm{def:b3:galois:separable}{séparable}, non réflexif. Signatures~:
$c_0$ --- bidual $\ell^\infty$~; $\ell^1$ --- Schur~;
$\ell^p$ --- réflexivité~; $\ell^\infty$ --- non-séparabilité
et limites de Banach.

\textbf{23.} Suite minimisante $(f_k)$ bornée~: sous-suite
faiblement convergente $f_{k_j} \rightharpoonup f$
(question~13). Si $f \notin F$, Hahn--Banach produit
$\Lambda$ s'annulant sur $F$ avec $\Lambda(f) = 1$~: $0 =
\Lambda(f_{k_j}) \to 1$, contradiction. Donc $f \in F$, et
par semi-continuité inférieure faible $\norm{x - f} = d$.
Dans $c_0$~: $\norm\Lambda = 1$ pour $\Lambda(x) =
\sum_n2^{-n}x_n$ n'est jamais atteint (une suite nulle non
nulle ne peut saturer $\abs{x_n} = \norm x_\infty$ pour tout
$n$)~; la formule de distance $\operatorname{dist}(x,
\ker\Lambda) = \abs{\Lambda(x)}$ implique qu'aucun point le
plus proche n'existe hors de $\ker\Lambda$.

\textbf{24.} Extraire $(y_j)$ avec $\abs{\langle y_i,
y_j\rangle} \leq \frac1j$ pour $i < j$ (chaque crochet tend
vers $0$ par convergence faible). Alors
$\norm{\frac1m\sum_j y_j}_2^2 \leq \frac{C^2+2}m \to 0$~:
moyennes de Cesàro convergentes en norme (Banach--Saks). Sur
$(e_k)$~: $\norm{\frac1m(e_1+\dots+e_m)}_2 =
\frac1{\sqrt m}$.

\textbf{25.} $L(A_mx) = L(x)$ par linéarité et invariance par
décalage. Les bornes $\liminf A_mx \leq L(x) \leq \limsup
A_mx$ suivent de la positivité. Sur les suites $T$-périodiques,
$A_Tx$ est la constante égale à la moyenne de période~: toutes
les limites de Banach coïncident. Pour la suite par blocs
(uns sur $\intoc{3^{j-1}}{3^j}$ pour $j$ pair)~: les
moyennes de Cesàro oscillent entre $\leq \frac13$ et $\geq
\frac23$. Étendre depuis $c \oplus \R x$ avec $\Lambda_\pm(y
+ tx) = \lim y \pm t\,p(\pm x)$ dominées par $p$~: deux
limites de Banach $L_\pm$ avec $L_-(x) \leq \frac13 <
\frac23 \leq L_+(x)$.


Compléments détaillés pour les parties~IV--VI.

Sur la convergence faible (questions~13--16)~: l'extraction
diagonale sur les coordonnées d'une suite bornée de $\ell^p$
($1 < p < \infty$) produit une limite coordonnée $x$ qui
appartient à $\ell^p$ par les bornes de sections finies
$\sum_{n\leq N}\abs{x_n}^p \leq M^p$. Pour vérifier la
convergence faible contre tout $y \in \ell^q$, découper le
crochet en une tête finie (qui tend vers $0$ coordonnée par
coordonnée) et une queue contrôlée par Hölder et la
bornitude uniforme des normes. Dans $\ell^1$ le même
raisonnement échoue~: $(e_n)$ est bornée, toute sous-suite
converge coordonnée par coordonnée vers $0$, mais le crochet
contre la suite de signes alternés $(-1)^k$ sur les indices
extraits diverge. Banach--Steinhaus appliqué aux $J
x^{(k)}$ sur le \omterm{def:b3:banach:operator}{dual} $E'$ (qui est toujours \omterm{def:b3:complete:complete}{complet}) convertit
la convergence des crochets en bornitude des normes~: toute
suite faiblement convergente est bornée.

Sur Schur (questions~19--21)~: la construction récursive des
blocs isole la masse de chaque $x^{(k_j)}$ sur un intervalle
d'indices $B_j$ disjoint des précédents, en tuant d'abord la
tête (convergence coordonnée) puis la queue (convergence de
la série définissant la norme $\ell^1$). Le signe de $y$ sur
chaque bloc récupère presque toute la norme, contredisant la
convergence faible vers $0$. Ainsi dans $\ell^1$ les
convergences faible et en norme des suites coïncident ---
pourtant les \omterm{def:b3:topology:topology}{topologies} diffèrent (les \omterm{def:b3:topology:topology}{voisinages} faibles ne
sont jamais bornés en norme) et la boule unité n'est pas
séquentiellement faiblement compacte, ce qui est cohérent avec
la non-réflexivité.

Sur les points les plus proches et Banach--Saks
(questions~23--24)~: la réflexivité fournit une sous-suite
faiblement convergente d'une suite minimisante~; Hahn--Banach
empêche la limite faible de sortir de $F$ (sinon une
fonctionnelle s'annulant sur $F$ séparerait)~; la
semi-continuité inférieure faible de la norme donne
l'atteinte. Dans $c_0$ la fonctionnelle $\Lambda(x) =
\sum_n 2^{-n}x_n$ n'atteint pas sa norme sur la boule unité
(une suite nulle ne peut saturer $\abs{x_n} = \norm
x_\infty$ pour tout $n$), et la formule de distance
$\operatorname{dist}(x, \ker\Lambda) = \abs{\Lambda(x)}$
implique qu'aucun point hors de $\ker\Lambda$ n'a de
projection sur cet hyperplan fermé. Banach--Saks dans
$\ell^2$ extrait une sous-suite presque orthogonale (chaque
crochet croisé $< \frac1j$) dont les moyennes de Cesàro
tendent vers $0$ en norme, convertissant la convergence
faible en convergence en norme pour les moyennes --- visible
sur $(e_k)$ où $\norm{\frac1m\sum_{j=1}^m e_j}_2 =
\frac1{\sqrt m}$.

Sur les limites de Banach (question~25)~: $L(A_mx) = L(x)$
par linéarité et invariance par décalage, et les bornes
$\liminf A_mx \leq L(x) \leq \limsup A_mx$ suivent de la
positivité appliquée aux queues. Sur les suites périodiques,
$A_Tx$ est la suite constante égale à la moyenne de période~:
toutes les limites de Banach coïncident. Pour la suite par
blocs égale à $1$ sur $\intoc{3^{j-1}}{3^j}$ pour $j$ pair,
les moyennes de Cesàro en $N = 3^j$ oscillent entre $\leq
\frac13$ ($j$ impair) et $\geq \frac23$ ($j$ pair).
L'extension depuis $c \oplus \R x$ avec les valeurs
extrêmes $\pm p(\pm x)$ est dominée par le sous-linéaire
$p$ (sous-additivité sur $y \pm sx$), et Hahn--Banach
produit deux limites de Banach $L_\pm$ séparant $x$~: hors
du monde périodique (et plus généralement
presque-convergent), une limite de Banach est un véritable
choix.


Détails supplémentaires pour les parties~I--III. L'inégalité
de Young $ab \leq \frac{a^p}p + \frac{b^q}q$ se lit comme
la concavité de $\log$ (ou comme le minimum de $t \mapsto
\frac{t^p}p + \frac1q - t$ en $t = 1$ quand $b = 1$).
Hölder s'en déduit par sommation après normalisation, et
Minkowski en écrivant $\abs{x_n+y_n}^p \leq
(\abs{x_n}+\abs{y_n})\abs{x_n+y_n}^{p-1}$ puis en appliquant
Hölder avec l'exposant conjugué (noter $(p-1)q = p$). La
dualité $(\ell^p)' = \ell^q$ repose sur le test
$x_n = \abs{y_n}^{q-1}\operatorname{sign}(y_n)$ (tronqué)~:
le crochet récupère $\sum\abs{y_n}^q$ tandis que la norme
$\ell^p$ est $(\sum\abs{y_n}^q)^{1/p}$, d'où
$\norm{\Lambda_y} = \norm y_q$ après passage à la limite
sur la longueur de troncature. La réciproque pose $y_n =
\Lambda(e_n)$ et utilise les mêmes tests pour obtenir $y
\in \ell^q$~; l'accord sur les suites finies \omterm{def:b3:topology:interior}{denses} donne
$\Lambda = \Lambda_y$. La réflexivité suit en composant les
deux dualités et en vérifiant que le composé est le $J$
canonique. La non-séparabilité de $\ell^\infty$ (famille non
dénombrable de suites indicatrices à distance mutuelle $1$)
bloque la réflexivité de $\ell^1$ via
\cref{exo:b3:banach:10}. Les limites de Banach, construites
par Hahn--Banach depuis le sous-linéaire $p(x) =
\limsup\frac{x_1+\dots+x_n}{n}$, vivent hors de $\ell^1$
(elles s'annulent sur les $e_n$ mais pas sur $\mathbf 1$) et
illustrent concrètement $(\ell^\infty)' \supsetneq
\ell^1$.

\end{solution}
